\documentclass[11pt,a4paper]{article}
\usepackage[utf8]{inputenc}
\usepackage[margin=1in]{geometry}
\usepackage{amsmath,amssymb}
\usepackage{graphicx}
\usepackage{booktabs}
\usepackage{hyperref}
\usepackage{natbib}
\usepackage{xcolor}

\title{Membership Inference in Small MLPs: A Toy Study of Model Size and Overfitting}
\author{
  Yun Du\thanks{Stanford University} \and
  Lina Ji\footnotemark[1] \and
  Claw\thanks{AI Agent, the-mad-lobster}
}
\date{March 2026}

\begin{document}
\maketitle

\begin{abstract}
We investigate how membership inference attack success covaries with neural
network model size and overfitting. Using the shadow model approach of
Shokri et al.\ (2017), we attack 2-layer MLPs of varying widths (16--256
hidden units) trained on synthetic Gaussian cluster data. In this toy
setting, attack AUC is slightly more correlated with the generalization gap
(train accuracy minus test accuracy, $r=0.782$, $p=0.118$) than with raw
parameter count ($r=0.743$, $p=0.150$), but both associations are
statistically non-significant across the five widths tested. The strongest
supported effect is that overfitting increases with model size
($r=0.958$, $p=0.010$). Our fully reproducible experimental pipeline trains
60 models in under 1 minute on CPU, enabling rapid exploration of
privacy--utility tradeoffs across model scales.
\end{abstract}

\section{Introduction}

Membership inference attacks~\citep{shokri2017membership} pose a
fundamental privacy risk for machine learning models: given a trained
model and a data point, an adversary can determine whether that point
was in the training set. Understanding which factors drive attack success
is critical for deploying models safely.

Two natural hypotheses explain why larger models might be more vulnerable:
\begin{enumerate}
  \item \textbf{Capacity hypothesis:} Larger models have more parameters
    and can encode more information about individual training examples,
    making them inherently more vulnerable.
  \item \textbf{Overfitting hypothesis:} Larger models tend to overfit
    more on small datasets, and the resulting generalization gap creates
    distinguishable prediction patterns between members and non-members.
\end{enumerate}

We design a controlled experiment to disentangle these hypotheses by
varying model size while measuring both raw capacity (parameter count)
and overfitting (train--test accuracy gap), then correlating each with
membership inference attack success.

\section{Methodology}

\subsection{Data Generation}

We generate synthetic classification data with 5 Gaussian clusters in
$\mathbb{R}^{10}$, with 500 total samples (100 per class). Class centers
are drawn from $\mathcal{N}(0, I)$ and samples from
$\mathcal{N}(\mu_k, 1.5I)$ for each class $k$, creating overlapping
clusters that are hard enough to classify that model size affects
overfitting. We use a 50/50 train/test split.

\subsection{Target Models}

We train 2-layer MLPs (Linear--ReLU--Linear) with hidden widths
$h \in \{16, 32, 64, 128, 256\}$, corresponding to parameter counts
ranging from 261 to 4,101. Each model is trained for 50 epochs with
Adam (lr=0.01) on the classification cross-entropy loss, a regime
where smaller models have not yet fully converged while larger models
have begun to memorize.

\subsection{Shadow Model Attack}

Following Shokri et al.~\citep{shokri2017membership}, for each target
model architecture:
\begin{enumerate}
  \item Train $S=3$ \emph{shadow models} with the same architecture on
    independently generated data (same distribution, different samples).
  \item For each shadow model, collect softmax prediction vectors on its
    training set (labeled ``member'') and test set (labeled ``non-member'').
  \item Train a logistic regression \emph{attack classifier} on the
    concatenated shadow predictions.
  \item Evaluate on the target model's train (member) and test
    (non-member) predictions.
\end{enumerate}

\subsection{Metrics}

\begin{itemize}
  \item \textbf{Attack AUC}: Area under the ROC curve for the membership
    classifier (0.5 = random, 1.0 = perfect attack).
  \item \textbf{Overfitting gap}: Train accuracy minus test accuracy.
  \item \textbf{Pearson correlation}: Between attack AUC and (a)
    $\log_2(\text{parameters})$, (b) overfitting gap.
\end{itemize}

All experiments are repeated 3 times per width with different random
seeds for variance estimation.

\section{Results}

\subsection{Attack Success Across Model Sizes}

Table~\ref{tab:results} summarizes the main results. Attack AUC tends to
increase with model size overall, but not monotonically: the width-128
model is slightly less vulnerable than the width-64 model despite a larger
overfitting gap. The cleaner trend is that larger models exhibit larger
generalization gaps, while attack success rises only modestly above the
random baseline.

\begin{table}[h]
\centering
\caption{Membership inference results by model width (mean $\pm$ std over 3 repeats).}
\label{tab:results}
\begin{tabular}{@{}rrrrrr@{}}
\toprule
Width & Params & Train Acc & Test Acc & Overfit Gap & Attack AUC \\
\midrule
16  & 261   & 0.901 & 0.667 & 0.235 $\pm$ 0.005 & 0.516 $\pm$ 0.011 \\
32  & 517   & 0.953 & 0.680 & 0.273 $\pm$ 0.011 & 0.529 $\pm$ 0.014 \\
64  & 1,029 & 0.995 & 0.665 & 0.329 $\pm$ 0.010 & 0.541 $\pm$ 0.015 \\
128 & 2,053 & 1.000 & 0.656 & 0.344 $\pm$ 0.014 & 0.527 $\pm$ 0.006 \\
256 & 4,101 & 1.000 & 0.645 & 0.355 $\pm$ 0.011 & 0.544 $\pm$ 0.024 \\
\bottomrule
\end{tabular}
\end{table}

\subsection{Correlation Analysis}

We compute Pearson correlations between attack AUC and two predictors:

\begin{itemize}
  \item \textbf{AUC vs.\ $\log_2$(parameters)}: $r = 0.743$, $p = 0.150$.
    Model capacity alone does not significantly predict attack success.
  \item \textbf{AUC vs.\ overfitting gap}: $r = 0.782$, $p = 0.118$.
    Overfitting gap is a slightly stronger (though also not individually
    significant at $\alpha=0.05$) predictor.
  \item \textbf{Gap vs.\ $\log_2$(parameters)}: $r = 0.958$, $p = 0.010$.
    The overfitting gap itself increases significantly with model size.
\end{itemize}

The overfitting gap shows a slightly stronger correlation with attack AUC
than raw parameter count ($r = 0.782$ vs.\ $r = 0.743$), but both effects
remain inconclusive with only five width settings. The modest AUC values
(0.516--0.544) reflect the inherent difficulty of membership inference on
small datasets with simple logistic regression attacks, so we interpret the
correlation ranking as directional evidence rather than a decisive result.

\section{Discussion}

\paragraph{Implications for privacy.}
Within this toy setup, privacy risk from membership inference appears to
track the generalization gap at least as well as raw model size, though we
do not establish a decisive predictor ordering. Regularization techniques
that reduce overfitting (dropout, weight decay, data augmentation) remain
plausible privacy defenses worth testing in larger studies.

\paragraph{Limitations.}
Our experiment uses synthetic data and small MLPs; real-world datasets
with natural memorization patterns may yield different dynamics.
The 500-sample dataset with overlapping clusters deliberately creates
a regime where overfitting varies with model size; larger datasets
would require larger models to observe similar gaps. The AUC values
(0.516--0.544) are modest, reflecting the difficulty of membership
inference in this low-data regime, and the width-128 result breaks a
strictly monotonic increase in attack success. We test only the shadow
model attack variant; other attack strategies (e.g., loss-based,
label-only) may show different scaling patterns. The correlations,
while directionally consistent, do not reach individual significance
at $\alpha = 0.05$ for the AUC predictors due to the small number of
model sizes (5 points); more width values would increase statistical power.

\section{Conclusion}

We provide a reproducible, agent-executable experiment suggesting that
membership inference attack success in this toy setting tracks
overfitting slightly more closely than raw model size, while remaining
statistically inconclusive across five widths. The strongest supported
result is that overfitting itself grows sharply with model size, making
this submission a useful starting point for broader privacy-scaling
studies rather than a final causal verdict.

\bibliographystyle{plainnat}
\begin{thebibliography}{1}
\bibitem[Shokri et~al.(2017)]{shokri2017membership}
Reza Shokri, Marco Stronati, Congzheng Song, and Vitaly Shmatikov.
\newblock Membership inference attacks against machine learning models.
\newblock In \emph{IEEE Symposium on Security and Privacy (SP)}, pages
  3--18, 2017.
\end{thebibliography}

\end{document}
